\documentclass[12pt,a4paper]{article}

\usepackage[margin=1in]{geometry}
\usepackage[T1]{fontenc}
\usepackage[utf8]{inputenc}
\usepackage{array}
\usepackage{booktabs}
\usepackage{float}
\usepackage{xcolor}
\usepackage{listings}
\usepackage{graphicx}

\lstset{
    basicstyle=\ttfamily\footnotesize,
    backgroundcolor=\color{gray!10},
    breaklines=true,
    columns=fullflexible,
    frame=single,
    framesep=6pt,
    rulecolor=\color{gray!55},
    numbers=left,
    numberstyle=\tiny,
    keywordstyle=\color{blue},
    commentstyle=\color{gray},
    stringstyle=\color{teal},
    showstringspaces=false,
    tabsize=4
}

\title{CPU Scheduling Algorithms Lab Report}
\author{Name: Abu Bakar Siddique\\Roll: 47}
\date{\today}

\begin{document}

\maketitle

\section{Theory}

CPU scheduling decides which ready process will use the CPU next. In this lab, the following algorithms were implemented.

\subsection{First-Come, First-Served}

FCFS is a non-preemptive algorithm. Processes execute in the same order in which they arrive. For processes with arrival time 0:

\[
WT_i = \sum_{j=1}^{i-1} BT_j
\]

\[
TAT_i = WT_i + BT_i
\]

\subsection{Shortest Job First}

SJF is a non-preemptive algorithm. It selects the process with the smallest CPU burst time. If two processes have the same burst time, FCFS order is preserved.

\subsection{Shortest Remaining Time Next}

SRTN is the preemptive version of SJF. At each time unit, the process with the minimum remaining burst time is selected.

\[
TAT_i = CT_i - AT_i
\]

\[
WT_i = TAT_i - BT_i
\]

\subsection{Round Robin}

Round Robin is a preemptive scheduling algorithm. Each process gets the CPU for at most one time quantum. If it is not finished, it goes back to the ready queue.

\subsection{Lottery Scheduling}

Lottery Scheduling is probabilistic. Each process receives a number of tickets. At each scheduling decision, a random ticket is drawn, and the process holding that ticket gets the CPU for one time quantum.

\newcommand{\resultimage}[2]{
\begin{figure}[H]
\centering
\includegraphics[width=0.95\textwidth,height=0.45\textheight,keepaspectratio]{#1}
\caption{#2}
\end{figure}
}

\section{Results}

\subsection{FCFS Result}

Input burst times: 6, 8, 3, 7, 4 ms.

\begin{table}[H]
\centering
\caption{FCFS scheduling output}
\begin{tabular}{cccc}
\toprule
\textbf{Process} & \textbf{Burst Time} & \textbf{Waiting Time} & \textbf{Turnaround Time} \\
\midrule
P1 & 6 & 0 & 6 \\
P2 & 8 & 6 & 14 \\
P3 & 3 & 14 & 17 \\
P4 & 7 & 17 & 24 \\
P5 & 4 & 24 & 28 \\
\bottomrule
\end{tabular}
\end{table}

\resultimage{fcfs.png}{FCFS output screenshot}

Average waiting time = 12.20 ms. Average turnaround time = 17.80 ms.

\subsection{SJF Result}

The same process set was used for SJF.

\begin{table}[H]
\centering
\caption{SJF scheduling output}
\begin{tabular}{cccc}
\toprule
\textbf{Process} & \textbf{Burst Time} & \textbf{Waiting Time} & \textbf{Turnaround Time} \\
\midrule
P3 & 3 & 0 & 3 \\
P5 & 4 & 3 & 7 \\
P1 & 6 & 7 & 13 \\
P4 & 7 & 13 & 20 \\
P2 & 8 & 20 & 28 \\
\bottomrule
\end{tabular}
\end{table}

\resultimage{sjf.png}{SJF output screenshot}

Average waiting time = 8.60 ms. Average turnaround time = 14.20 ms.

\subsection{FCFS and SJF Comparison}

\begin{table}[H]
\centering
\caption{Comparison between FCFS and SJF}
\begin{tabular}{ccc}
\toprule
\textbf{Algorithm} & \textbf{Average Waiting Time} & \textbf{Average Turnaround Time} \\
\midrule
FCFS & 12.20 ms & 17.80 ms \\
SJF & 8.60 ms & 14.20 ms \\
\bottomrule
\end{tabular}
\end{table}

For the same process set, SJF gives a lower average waiting time than FCFS because shorter jobs are executed earlier.

\subsection{SRTN Result}

\begin{table}[H]
\centering
\caption{SRTN scheduling output}
\begin{tabular}{cccccc}
\toprule
\textbf{Process} & \textbf{Arrival} & \textbf{Burst} & \textbf{Completion} & \textbf{Waiting} & \textbf{Turnaround} \\
\midrule
P1 & 0 & 8 & 17 & 9 & 17 \\
P2 & 1 & 4 & 5 & 0 & 4 \\
P3 & 2 & 9 & 26 & 15 & 24 \\
P4 & 3 & 5 & 10 & 2 & 7 \\
\bottomrule
\end{tabular}
\end{table}

\resultimage{srtn.png}{SRTN output screenshot}

Average waiting time = 6.50 ms. Average turnaround time = 13.00 ms. This demonstrates preemption because P1 starts first, but P2 arrives at time 1 with a shorter remaining time and preempts P1.

\subsection{Round Robin Result}

Input burst times: 24, 3, 3 ms.

\begin{table}[H]
\centering
\caption{Round Robin comparison}
\begin{tabular}{ccc}
\toprule
\textbf{Quantum} & \textbf{Average Waiting Time} & \textbf{Average Turnaround Time} \\
\midrule
4 ms & 5.67 ms & 15.67 ms \\
2 ms & 6.33 ms & 16.33 ms \\
\bottomrule
\end{tabular}
\end{table}

\resultimage{rr.png}{Round Robin output screenshot}

For this input, quantum 4 ms performs slightly better than quantum 2 ms in both average waiting time and average turnaround time.

\subsection{Lottery Scheduling Result}

Input: P1 burst = 10, tickets = 10; P2 burst = 6, tickets = 20; P3 burst = 4, tickets = 5; quantum = 2 ms.

\begin{table}[H]
\centering
\caption{Lottery Scheduling three-run comparison}
\begin{tabular}{ccc}
\toprule
\textbf{Run} & \textbf{Average Waiting Time} & \textbf{Average Turnaround Time} \\
\midrule
1 & 8.67 ms & 15.33 ms \\
2 & 10.00 ms & 16.67 ms \\
3 & 8.67 ms & 15.33 ms \\
\bottomrule
\end{tabular}
\end{table}

\resultimage{lottery.png}{Lottery Scheduling output screenshot}

The three runs produced different execution orders, which shows the probabilistic nature of Lottery Scheduling.

\appendix

\section{Source Code}

\lstinputlisting[language=C,caption={scheduling.c}]{scheduling.c}

\end{document}
